\documentclass[12pt,a4paper]{article}

% Paket
\usepackage[utf8]{inputenc}
\usepackage[T1]{fontenc}
\usepackage[swedish]{babel}
\usepackage[autostyle=true]{csquotes}
\usepackage[margin=2.5cm]{geometry}
\usepackage{graphicx}
\usepackage{xcolor}
\usepackage{hyperref}
\usepackage{amsmath}
\usepackage{listings}
\usepackage{upquote}
\usepackage{tikz}
\usepackage{enumitem}
\usepackage{booktabs}
\usepackage{array}
\usepackage{verbatim}

% Undvik konflikt med babels "-shorthand i löpande text
\AtBeginDocument{\shorthandoff{"}}

% Färger
\definecolor{brandgreen}{RGB}{0,112,74}
\definecolor{lightgray}{RGB}{245,245,245}
\definecolor{codegray}{RGB}{50,50,50}

% Hyperref inställningar
\hypersetup{
  colorlinks=true,
  linkcolor=brandgreen,
  urlcolor=brandgreen,
  citecolor=brandgreen
}

% Listings för kodexempel
\lstset{
  basicstyle=\small\ttfamily,
  backgroundcolor=\color{lightgray},
  breaklines=true,
  frame=single,
  numbers=left,
  numberstyle=\tiny\color{codegray},
  literate=%
    {å}{{\aa}}1
    {ä}{{\"a}}1
    {ö}{{\"o}}1
    {Å}{{\AA}}1
    {Ä}{{\"A}}1
    {Ö}{{\"O}}1
    {é}{{\'e}}1
    {≥}{$\geq$}1
    {->}{{$\to$}}2
    {→}{{$\to$}}1
    {!=}{{$\neq$}}2
    {≠}{{$\neq$}}1
}

% Unicode-mappningar för pdfLaTeX (vanlig text)
\DeclareUnicodeCharacter{2192}{$\rightarrow$} % →
\DeclareUnicodeCharacter{2260}{$\neq$}       % ≠
\DeclareUnicodeCharacter{2705}{}              % ✅ (ta bort)

% Definiera JSON och JavaScript för listings
\lstdefinelanguage{json}{
  morestring=[b]",
  morestring=[d]'
}

\lstdefinelanguage{javascript}{
  keywords={const, let, var, function, if, else, return, await, async, try, catch, fetch},
  morecomment=[l]{//},
  morestring=[b]",
  morestring=[d]'
}

% Titel
\title{\textbf{API Endpoints Specifikation: Content Slides} \\ \large Backend-mappning för Onboarding-flödet (med Sidebar)}
\author{Celestial AB}
\date{\today}

\begin{document}

\maketitle
\tableofcontents
\newpage

% ========================================================================
\section{Introduktion}
% ========================================================================

Detta dokument beskriver \textbf{alla API endpoints} som behövs för onboarding-slidesen (content slides med sidebar) enligt \texttt{Onboardin\_app\_ny.tex}.

Dessa slides visas \textbf{efter inloggning} och utgör själva onboarding-processen där användaren:
\begin{itemize}[noitemsep]
  \item Väljer uppdragets art (vilka tjänster som behövs)
  \item Besvarar riskfrågor (PTL-krav)
  \item Genomgår identitetskontroll
  \item Jämförs mot Bolagsverket/Roaring.io
  \item Laddar upp företagsdokumentation
  \item Får ekonomisk analys och riskbedömning
  \item Signerar avtal med BankID
\end{itemize}

\subsection{Dokumentstruktur}

Dokumentet följer samma struktur som \texttt{Onboardin\_app\_ny.tex} med sections och subsections:

\begin{itemize}[noitemsep]
  \item \textbf{Sektion 1:} Uppdragsval och introduktion
  \item \textbf{Sektion 2:} Riskbedömning (Steg 1 + Steg 2 fördjupning)
  \item \textbf{Sektion 3:} Identitetskontroll
  \item \textbf{Sektion 4:} Jämförelse mot externa register
  \item \textbf{Sektion 5:} Företagsdokumentation
  \item \textbf{Sektion 6:} Ekonomisk analys (SIE + Bankkonto)
  \item \textbf{Sektion 7:} Djupgranskning och avtal
\end{itemize}

För varje slide/steg dokumenteras:
\begin{itemize}[noitemsep]
  \item Vilka API endpoints som behövs
  \item HTTP-metod (GET, POST, PUT, DELETE)
  \item Request-format (JSON schema, multipart/form-data för filer)
  \item Response-format (success och error cases)
  \item Affärslogik och valideringar
  \item Externa integrationer (Bolagsverket, Skatteverket, Roaring.io, BankID, etc.)
  \item Databastabeller som påverkas
  \item Riskindikatorer och flaggor
\end{itemize}

\subsection{Teknisk arkitektur}

\begin{itemize}[noitemsep]
  \item \textbf{Databas:} CSV-filer under utveckling, PostgreSQL som slutoptimering (använder UUID för process-ID)
  \item \textbf{Externa API:er:} Bolagsverket, Skatteverket SPAR, Roaring.io, BankID
  \item \textbf{AI-analys:} Hugging Face (microsoft/layoutlmv3-large) via Celery background tasks
  \item \textbf{Fillagring:} S3/MinIO för uppladdade dokument (SIE, PDF, etc.)
  \item \textbf{Autentisering:} JWT tokens med user\_id, varje onboarding-process identifieras med UUID
\end{itemize}

\textbf{VIKTIGT om databasval:} Hela applikationen byggs FÖRST med CSV-filer för ALL databashantering (onboarding\_processes.csv, external\_api\_responses.csv, event\_log.csv, etc.). PostgreSQL SQL införs som slutoptimering ENDAST efter att hela applikationen fungerar felfritt. Detta för att undvika race conditions och skrivfel vid flera samtidiga användare. CSV-first approach = snabbare utveckling, enklare debugging.

\newpage

% ========================================================================
\section{SEKTION 1: Uppdragsval och introduktion}
% ========================================================================

\subsection{Content Slide: Välkommen – Uppdragets art}

\textbf{React route:} \texttt{/inledning}\\
\textbf{Beamer label:} \texttt{intro}\\
\textbf{Sidebar:} Ja (första slide med sidebar)

\textbf{Beskrivning:}\\
Användaren väljer vilka tjänster de är intresserade av (löpande bokföring, årsbokslut, deklarationer, etc.). Detta är första steget i onboarding och sätter tonen för hela processen – kunden i centrum, myndighetskrav sekundärt.

% ------------------------------------------------------------------------
\subsection{API-endpoints för Uppdragsval och introduktion}
% ------------------------------------------------------------------------

\subsubsection{POST /api/onboarding/uppdrag}

\textbf{Beskrivning:}\\
Skapar en ny onboarding-process, sparar kundens valda tjänster och beräknar uppskattad kostnad baserat på vald tjänstekombination. Returnerar ett UUID som används för alla efterföljande API-anrop i denna onboarding-process.

\textbf{Request Headers:}
\begin{lstlisting}[language=json]
{
  "Authorization": "Bearer <access_token>",
  "Content-Type": "application/json"
}
\end{lstlisting}

\textbf{Request Body:}
\begin{lstlisting}[language=json]
{
  "lopandeBokforing": true,
  "arsbokslut": true,
  "deklarationer": false,
  "loneadministration": false,
  "ekonomiskRadgivning": true,
  "foretagsregistrering": false,
  "finansiellRapportering": false,
  "foretagsforsaljning": false,
  "annat": "Succession till nästa generation"
}
\end{lstlisting}

\textbf{Response (200 OK):}
\begin{lstlisting}[language=json]
{
  "success": true,
  "onboardingId": "a7f3c8e2-4b1d-4e9a-8c3f-1a2b3c4d5e6f",
  "uppskattadKostnad": "15000 kr/år",
  "antalTjanster": 3,
  "riskIndikator": "low",
  "message": "Onboarding-process skapad",
  "nextStep": "/riskfragor"
}
\end{lstlisting}

\textbf{Response (400 Bad Request):}
\begin{lstlisting}[language=json]
{
  "success": false,
  "error": "Minst en tjänst måste väljas",
  "code": "NO_SERVICES_SELECTED"
}
\end{lstlisting}

\textbf{Response (401 Unauthorized):}
\begin{lstlisting}[language=json]
{
  "success": false,
  "error": "Ogiltig eller utgången token",
  "code": "INVALID_TOKEN"
}
\end{lstlisting}

\textbf{Business Logic:}
\begin{enumerate}[noitemsep]
  \item Hämta \texttt{user\_id} från JWT token (Authorization header)
  \item Validera att minst en tjänst är vald
  \item Generera nytt UUID för onboarding-processen
  \item Beräkna uppskattad kostnad baserat på prislista:
    \begin{itemize}[noitemsep]
      \item Löpande bokföring: 5000 kr/år
      \item Årsbokslut: 8000 kr/år
      \item Deklarationer: 3000 kr/år
      \item Löneadministration: 4000 kr/år
      \item Ekonomisk rådgivning: 10000 kr/år
      \item Företagsregistrering: 15000 kr (engångsavgift)
      \item Finansiell rapportering: 6000 kr/år
      \item Företagsförsäljning/succession: Offert (kontakt)
    \end{itemize}
  \item Sätt riskindikator:
    \begin{itemize}[noitemsep]
      \item \texttt{foretagsregistrering = true} → riskIndikator = \enquote{medium} (nystartad = högre PTL-risk)
      \item Annars → riskIndikator = \enquote{low}
    \end{itemize}
  \item Spara i databas (CSV under utveckling, PostgreSQL som slutoptimering)
  \item Returnera \texttt{onboardingId} (UUID), uppskattad kostnad och nästa steg
\end{enumerate}

\textbf{Databas (SQL-syntax för framtida PostgreSQL, implementeras först som CSV):}
\begin{lstlisting}[language=sql]
-- Skapa ny onboarding-process
-- NUVARANDE IMPLEMENTATION: onboarding_processes.csv
-- FRAMTIDA OPTIMERING: PostgreSQL (efter att hela appen fungerar)
INSERT INTO onboarding_processes (
  id,                    -- UUID genererat av backend
  user_id,               -- Från JWT token
  status,                -- 'draft'
  current_step,          -- 1
  uppdrag_data,          -- JSONB med valda tjänster
  uppskattad_kostnad,    -- Beräknad kostnad
  risk_indikator,        -- 'low' eller 'medium'
  created_at,
  updated_at
) VALUES (
  gen_random_uuid(),     -- PostgreSQL genererar UUID
  $1,                    -- user_id
  'draft',
  1,
  $2::jsonb,             -- uppdrag_data
  $3,                    -- uppskattad_kostnad
  $4,                    -- risk_indikator
  NOW(),
  NOW()
) RETURNING id;          -- Returnera UUID till backend
\end{lstlisting}

\textbf{Externa integrationer:}\\
Inga (ren datalagring)

\newpage
% ========================================================================
\section{Sektion 2A: Riskfrågor – Steg 1}
% ========================================================================

% ------------------------------------------------------------------------
\subsection{API-endpoints Sammanfattning (Våra endpoints)}
% ------------------------------------------------------------------------

\begin{itemize}[noitemsep]
  \item \texttt{POST /api/onboarding/\{onboardingId\}/riskfragor/steg1} — Sparar grundfrågor, validerar mot externa API:er, beräknar triggers
  \item \textbf{Antal anrop vid sidladdning:} 0 (endpointen anropas när användaren klickar \enquote{Nästa})
\end{itemize}

% ------------------------------------------------------------------------
\subsection{Översikt - När anropas våra endpoints}
% ------------------------------------------------------------------------

Riskfrågor Steg 1 är första delen i en multi-stage wizard för riskbedömning enligt Penningtvättslagen (PTL). Användaren presenteras med 9 grundläggande frågor om företagets verksamhet, kunder och ägare.

\textbf{React route:} \texttt{/riskfragor/steg1}\\
\textbf{React komponent:} \texttt{RiskFragorSteg1.jsx}\\
\textbf{Beamer slide:} \texttt{riskfragor-steg1}

\textbf{De 9 grundfrågorna:}
\begin{enumerate}[noitemsep]
  \item Företagets huvudsakliga affärsidé (textarea)
  \item Kundtyper (checkboxes: privatpersoner, företag, offentlig sektor)
  \item Utländska affärspartners/kunder (textarea)
  \item Största leverantörer (textarea)
  \item Verksamhet ändrad senaste 12 månaderna (textarea)
  \item Organisationsnummer (input)
  \item Personnummer VD/Firmatecknare (input)
  \item \textbf{Är det du som är verklig huvudman?} (radio: Ja >25\%, Nej, Ej tillämpligt)
  \item PEP-status (checkbox)
\end{enumerate}

När användaren fyller i formuläret och klickar \enquote{Nästa} skickas ett POST-anrop till backend med alla svar. Backend validerar organisationsnummer och personnummer mot externa API:er (Bolagsverket, SPAR, Roaring), beräknar triggers för riskflaggor, och returnerar ALLTID:
\begin{itemize}[noitemsep]
  \item \texttt{nextStep = "/riskfragor/steg2"} (Steg 2-4 visas ALLTID, eller fler enligt config.json)
  \item \texttt{triggers} objekt används för risk-scoring och UI-betoning, INTE för att hoppa över steg
  \item Triggers kan påverka vilka frågor som betonas med varningar eller extra hjälptext
\end{itemize}

\textbf{Vad används triggers till?}
\begin{itemize}[noitemsep]
  \item Flagga risker i backend för compliance-team
  \item Betona specifika frågor i UI (t.ex. markera utländska transaktioner med varningsikon)
  \item Påverka slutlig risk-score (low/medium/high)
  \item Dokumentera varför Enhanced Due Diligence (EDD) kan krävas
\end{itemize}

% ------------------------------------------------------------------------
\subsection{Externa API-endpoints Sammanfattning}
% ------------------------------------------------------------------------

\begin{itemize}[noitemsep]
  \item \texttt{POST /organisationer} (Bolagsverket) — Validera org.nr, hämta VD/firmatecknare
  \item \texttt{POST /navet/spar/personnummer} (Skatteverket SPAR) — Validera personnummer, folkbokföring
  \item \texttt{GET /se/company/beneficial-owners/2.1/\{orgNr\}} (Roaring) — Hämta verkliga huvudmän >25\%
  \item \texttt{GET /se/company/alternative-beneficial-owner/1.0/\{orgNr\}} (Roaring) — Hämta alternativ VH (VD) om ingen >25\%
  \item \texttt{GET /se/pep/1.0/\{personnr\}} (Roaring) — PEP-kontroll per person
  \item \texttt{GET /se/sanctions/1.0/\{personnr\}} (Roaring) — Sanktionskontroll per person
  \item \texttt{GET /se/businessprohibition/1.0/person/\{personnr\}} (Roaring) — Näringsförbud per person
\end{itemize}

\textbf{Totalt antal externa anrop:} 3-9 beroende på antal verkliga huvudmän (se Implementation Notes)

% ------------------------------------------------------------------------
\subsection{Översikt - Varför behövs externa endpoints}
% ------------------------------------------------------------------------

PTL 3 kap 2 § kräver att vi identifierar och verifierar kunden (företaget) samt verkliga huvudmän. Detta innebär:

\begin{itemize}[noitemsep]
  \item \textbf{Bolagsverket:} Verifiera att organisationsnumret finns, hämta officiellt företagsnamn, SNI-kod, VD och firmatecknare. Endast personer registrerade som VD/firmatecknare får påbörja onboarding.
  \item \textbf{SPAR:} Verifiera att personnumret är korrekt och att personen är folkbokförd i Sverige. Hämta officiellt namn och adress för identitetskontroll.
  \item \textbf{Roaring Beneficial Owners:} PTL 3 kap 6 § kräver identifiering av verkliga huvudmän (fysiska personer som äger/kontrollerar >25\%). Om ingen når tröskeln används PTL 3 kap 8 § (alternativ VH = högsta befattningshavare/VD).
  \item \textbf{Roaring PEP/Sanctions/Business Prohibition:} PTL 3 kap 10 § kräver screening mot PEP-listor och sanktionslistor. Näringsförbud är en extra kontroll enligt bokföringslagen.
\end{itemize}

Alla externa API-svar sparas i \texttt{external\_api\_responses} för audit trail och GDPR-compliance (rätt till tillgång).

% ------------------------------------------------------------------------
\subsection{Endpoint 1: POST /api/onboarding/\{onboardingId\}/riskfragor/steg1}
% ------------------------------------------------------------------------

\subsubsection{Beskrivning}

Sparar grundläggande riskfrågor (9 st.), validerar org.nr/personnr mot externa API:er (Bolagsverket, SPAR, Roaring), beräknar triggers för fördjupningsfrågor (Steg 2), och detekterar eventuell mismatch mellan användarens VH-påstående och Roaring.io:s register.

\subsubsection{Request}

\textbf{Headers:}
\begin{lstlisting}[language=json]
{
  "Authorization": "Bearer <access_token>",
  "Content-Type": "application/json"
}
\end{lstlisting}

\textbf{Body (pseudokod):}
\begin{lstlisting}
{
  "affarsIde": "string (textarea)",
  "kundTyper": {
    "privatpersoner": boolean,
    "foretag": boolean,
    "offentligSektor": boolean
  },
  "utlandskaPartners": "string (textarea)",
  "storaLeverantorer": "string (textarea)",
  "verksamhetAndrad": "string (textarea)",
  "organisationsnummer": "string (format: XXXXXX-XXXX)",
  "personnummer": "string (format: YYYYMMDD-XXXX)",
  "verkligHuvudman": "enum: ja_over_25 | nej | ej_tillampligt",
  "isPEP": boolean
}
\end{lstlisting}

\subsubsection{Response 200 OK - Inga riskflaggor}

\textbf{VIKTIGT:} Denna response följer \textbf{Celestial Risk Engine v3.0} (se \texttt{docs/Theory/metod\_riskbedömning\_kund\_v3.tex} för fullständig metodbeskrivning).

\begin{lstlisting}
{
  "success": true,
  "bolagsverketData": {
    "foretagsnamn": "string",
    "organisationsform": "string",
    "sniKod": "string",
    "sniBeskrivning": "string",
    "registreringsdatum": "ISO-8601 date",
    "status": "string"
  },
  "sparData": {
    "namn": "string",
    "personnummer": "string",
    "folkbokforingsadress": "string",
    "status": "string"
  },
  "riskAssessment": {
    "staticKycScore": 12,
    "adjustedScore": 32,
    "hotBonus": 20,
    "sniHotFactor": 3,
    "riskLevel": "low",
    "requiresEDD": false,
    "components": {
      "publicData": 4,
      "crossValidation": 8,
      "bookkeeping": 0
    },
    "triggers": {
      "isPEP": false,
      "hasSanctions": false,
      "hasBusinessProhibition": false,
      "vhMismatch": false,
      "addressMismatch": true,
      "missingFTax": false,
      "missingVAT": false,
      "newCompany": false,
      "highRiskIndustry": true
    }
  },
  "vhMismatch": {
    "detected": false,
    "reason": null,
    "requiresManualReview": false
  },
  "wizardConfig": {
    "totalSteps": 4,
    "hasCustomConfig": false,
    "steps": [
      {
        "id": 2,
        "title": "Utländska transaktioner",
        "isCustom": false
      },
      {
        "id": 3,
        "title": "Kunder & Affärspartners",
        "isCustom": false
      },
      {
        "id": 4,
        "title": "Betalningar & Transaktioner",
        "isCustom": false
      }
    ]
  },
  "message": "Grundläggande riskbedömning klar",
  "nextStep": "/riskfragor/steg2"
}
\end{lstlisting}

\textbf{OBS:} nextStep är ALLTID \enquote{/riskfragor/steg2}. Backend läser användarens custom config.json (om uppladdad via Inställningar) och mergear med default Steg 2-4. Frontend renderar wizard-steg dynamiskt baserat på \texttt{wizardConfig.steps}.

\textbf{Förklaring av riskAssessment:}
\begin{itemize}[noitemsep]
  \item \texttt{staticKycScore}: Summan av publicData + crossValidation + bookkeeping (bookkeeping = 0 i Steg 1)
  \item \texttt{hotBonus}: (sniHotFactor - 1) × 10 (Hot 1 = 0, Hot 2 = 10, Hot 3 = 20, Hot 4 = 30)
  \item \texttt{adjustedScore}: staticKycScore + hotBonus
  \item \texttt{sniHotFactor}: 1–4 från SNI-kod (se v3 metoddokument)
  \item \texttt{riskLevel}: low (0-40) / medium (41-70) / high (71-90) / critical (>90 eller hårda triggers)
  \item \texttt{requiresEDD}: true om Enhanced Due Diligence krävs (PTL 3 kap 10 §)
  \item \texttt{components}: Delpoäng per komponent (publicData max 30, crossValidation max 30, bookkeeping max 40)
  \item \texttt{triggers}: Boolean-flaggor för specifika riskfaktorer som påverkar score och åtgärder
\end{itemize}

\subsubsection{Response 200 OK - Riskflaggor aktiva}

\begin{lstlisting}
{
  "success": true,
  "bolagsverketData": { ... },
  "sparData": { ... },
  "triggers": {
    "showUtlandska": true,
    "showKontanter": true,
    "showVerksamhet": false,
    "showPEP": false
  },
  "vhMismatch": {
    "detected": false,
    "reason": null,
    "requiresManualReview": false
  },
  "wizardConfig": {
    "totalSteps": 4,
    "hasCustomConfig": false,
    "steps": [ /* samma som ovan */ ]
  },
  "message": "Riskflaggor identifierade (se triggers för detaljer)",
  "nextStep": "/riskfragor/steg2"
}
\end{lstlisting}

\textbf{OBS:} Samma nextStep och wizardConfig som vid inga triggers. Skillnaden är att triggers används för risk-scoring och kan påverka UI-presentation (betona frågor med varningsikon).

\subsubsection{Response 200 OK - VH Mismatch Detected (Warning)}

\begin{lstlisting}
{
  "success": true,
  "bolagsverketData": { ... },
  "sparData": { ... },
  "triggers": { ... },
  "vhMismatch": {
    "detected": true,
    "reason": "User claims VH but not found in Roaring register",
    "requiresManualReview": true,
    "details": {
      "userPersonnummer": "196505011234",
      "roaringVH": [
        {"personnummer": "198012251234", "shares": 45.5},
        {"personnummer": "197503201234", "shares": 30.0}
      ]
    }
  },
  "wizardConfig": {
    "totalSteps": 4,
    "hasCustomConfig": false,
    "steps": [ /* samma som ovan */ ]
  },
  "warning": "Din uppgift om verklig huvudman stämmer inte med offentliga register. En redovisningskonsult kommer att kontakta dig för klargörande.",
  "message": "Riskbedömning klar (manuell granskning krävs)",
  "nextStep": "/riskfragor/steg2"
}
\end{lstlisting}

\textbf{OBS:} VH-mismatch BLOCKERAR INTE processen. Användaren kan fortsätta, men ärendet flaggas för manuell granskning i Admin Dashboard.

\subsubsection{Response 200 OK - Med Custom Config}

\begin{lstlisting}
{
  "success": true,
  "bolagsverketData": { ... },
  "sparData": { ... },
  "triggers": { ... },
  "vhMismatch": { ... },
  "wizardConfig": {
    "totalSteps": 6,
    "hasCustomConfig": true,
    "customConfigVersion": "v2.1",
    "steps": [
      {
        "id": 2,
        "title": "Utländska transaktioner",
        "isCustom": false,
        "questionCount": 3
      },
      {
        "id": 3,
        "title": "Kunder & Affärspartners",
        "isCustom": false,
        "questionCount": 3
      },
      {
        "id": 4,
        "title": "Betalningar & Transaktioner",
        "isCustom": false,
        "questionCount": 5
      },
      {
        "id": 5,
        "title": "Miljöhänsyn & Hållbarhet",
        "isCustom": true,
        "questionCount": 4,
        "questions": [
          {
            "id": "env_1",
            "text": "Har företaget miljöcertifiering?",
            "type": "radio",
            "options": ["Ja", "Nej", "Under arbete"]
          },
          {
            "id": "env_2",
            "text": "Beskriv företagets miljöpolicy",
            "type": "textarea",
            "maxLength": 500
          }
          /* ... fler frågor ... */
        ]
      },
      {
        "id": 6,
        "title": "IT-säkerhet & GDPR",
        "isCustom": true,
        "questionCount": 6,
        "questions": [ /* ... */ ]
      }
    ]
  },
  "message": "Grundläggande riskbedömning klar (+ 2 anpassade steg)",
  "nextStep": "/riskfragor/steg2"
}
\end{lstlisting}

\textbf{OBS:} Backend läser användarens custom config.json (uppladdad via POST /api/user/settings/config) och mergear med default Steg 2-4. Custom steg kan INTE ersätta default steg 2-4, endast läggas till efter dem.

\subsubsection{Response 400 Bad Request}

\begin{lstlisting}
{
  "success": false,
  "error": "validation_error",
  "code": "INVALID_ORG_NUMBER",
  "details": "Organisationsnummer måste vara 10 siffror (XXXXXX-XXXX)"
}
\end{lstlisting}

\subsubsection{Response 403 Forbidden}

\begin{lstlisting}
{
  "success": false,
  "error": "not_authorized",
  "code": "PERSON_NOT_AUTHORIZED",
  "details": "Personnummer stämmer inte med VD/firmatecknare enligt Bolagsverket"
}
\end{lstlisting}

\subsubsection{Response 404 Not Found}
\begin{lstlisting}
{
  "success": false,
  "error": "not_found",
  "code": "COMPANY_NOT_FOUND",
  "orgNr": "556789-1234"
}
\end{lstlisting}

\subsubsection{Implementation Notes}

\textbf{Flöde (pseudokod):}
\begin{lstlisting}
1. Extrahera user_id från JWT token
2. SELECT onboarding_process WHERE id = onboardingId AND user_id = user_id
   OM INTE HITTAT -> 403 Forbidden
3. Validera request body format:
   - org.nr: Regex XXXXXX-XXXX
   - personnr: Regex YYYYMMDD-XXXX
   - verkligHuvudman: Enum ["ja_over_25", "nej", "ej_tillampligt"]
4. Anropa Bolagsverket API (POST /organisationer)
   SPARA response i external_api_responses
5. Anropa SPAR API (POST /navet/spar/personnummer)
   SPARA response i external_api_responses
6. Verifiera personnummer mot VD/firmatecknare från Bolagsverket
   OM INGEN MATCH -> 403 Forbidden
7. Anropa Roaring Beneficial Owners (två-stegs algoritm):
   a. GET /beneficial-owners/2.1/{orgNr}
   b. OM tom array -> GET /alternative-beneficial-owner/1.0/{orgNr}
   SPARA båda responses i external_api_responses
8. Kontrollera VH-mismatch (verklig huvudman-validering):
   OM verkligHuvudman === "ja_over_25":
     personnr_from_form = request.personnummer
     vh_from_roaring = external_api_responses.beneficial_owners
     OM personnr_from_form INTE finns i vh_from_roaring:
       FLAGGA mismatch_detected = true
       SPARA mismatch_reason = "User claims VH but not in Roaring register"
       LOG event: "VH_MISMATCH_DETECTED" med detaljer
       SET requires_manual_review = true
     ANNARS:
       mismatch_detected = false
9. FÖR varje VH (verklig huvudman från Roaring):
   a. GET /pep/1.0/{personnr}
   b. GET /sanctions/1.0/{personnr}
   c. GET /businessprohibition/1.0/person/{personnr}
   SPARA alla responses i external_api_responses
10. Beräkna triggers baserat på API-svar och avvikelser:
   
   A. BASERAD PÅ ANVÄNDARENS INPUT (sparas för riskbedömning):
   userFlags.hasInternationalPartners = (utlandskaPartners !== "Nej")
   userFlags.servesCash = (kundTyper.privatpersoner === true)
   userFlags.recentChanges = (verksamhetAndrad !== "Nej")
   userFlags.selfReportedPEP = (isPEP === true)
   
   B. TRIGGERS FRÅN API-AVVIKELSER (triggar extra åtgärder):
   triggers.pepDetected = (någon VH.isPEP === true FROM Roaring)
   triggers.newCompany = (registrationDate < 6 månader FROM Bolagsverket)
   triggers.addressMismatch = (SPAR adress !== Bolagsverket adress)
   triggers.missingFTax = (fTaxReg === false FROM Roaring)
   triggers.missingVAT = (vatReg === false FROM Roaring AND omsättning >80k)
   triggers.vhMismatch = (mismatch_detected === true FROM steg 8)
   triggers.sanctioned = (någon VH.isSanctioned === true FROM Roaring)
   triggers.businessProhibition = (någon VH.hasProhibition === true)
   triggers.highRiskIndustry = (SNI-kod i högrisklista FROM metoddokument)
   
   C. KOMBINATION (behövs för Celestial Engine):
   riskProfile = {
     userFlags: userFlags,
     apiTriggers: triggers,
     combinedRiskScore: calculate_celestial_score(userFlags, triggers)
   }

11. UPDATE onboarding_processes SET 
      riskfragor_steg1 = request_body,
      bolagsverket_data = bolagsverket_response,
      spar_data = spar_response,
      beneficial_owners_data = roaring_response,
      vh_mismatch_detected = mismatch_detected,
      vh_mismatch_reason = mismatch_reason,
      requires_manual_review = requires_manual_review,
      user_flags = userFlags,
      api_triggers = triggers,
      risk_profile = riskProfile,
      current_step = 2,
      updated_at = NOW()
    WHERE id = onboardingId
12. Läs användarens custom config (om uppladdad):
    custom_config = SELECT custom_wizard_config FROM user_settings WHERE user_id = user_id
    OM custom_config finns:
      wizard_config = merge_wizard_steps(DEFAULT_STEPS_2_4, custom_config.additionalSteps)
      wizard_config.hasCustomConfig = true
      wizard_config.totalSteps = len(wizard_config.steps)
    ANNARS:
      wizard_config = DEFAULT_STEPS_2_4
      wizard_config.hasCustomConfig = false
      wizard_config.totalSteps = 4
    OBS: Custom steg kan INTE ersätta Steg 2-4, endast läggas till som Steg 5, 6, 7...
13. Beräkna nextStep:
    nextStep = "/riskfragor/steg2" (ALLTID - Steg 2-N körs alltid)
    OBS: Triggers används för risk-scoring och UI-betoning, INTE för att hoppa över steg
14. Returnera response med nextStep="/riskfragor/steg2", wizardConfig, triggers (och mismatch_warning om flaggad)
\end{lstlisting}

\textbf{Valideringsregler:}
\begin{itemize}[noitemsep]
  \item org\_nr: Format XXXXXX-XXXX, måste finnas i Bolagsverket
  \item personnr: Format YYYYMMDD-XXXX, måste vara VD/firmatecknare enligt Bolagsverket
  \item personnr: Måste vara folkbokförd enligt SPAR
  \item verkligHuvudman: Enum ["ja\_over\_25", "nej", "ej\_tillampligt"]
  \item affarsIde: Max 1000 tecken
  \item utlandskaPartners: Max 1000 tecken
  \item storaLeverantorer: Max 1000 tecken
  \item verksamhetAndrad: Max 1000 tecken
\end{itemize}

\textbf{VH-Mismatch Hantering:}
\begin{itemize}[noitemsep]
  \item \textbf{Vad är en mismatch?} Användaren svarar "Ja, jag är VH >25\%" men deras personnummer finns INTE i Roaring.io:s VH-register
  \item \textbf{Varför kan detta hända?}
    \begin{itemize}[noitemsep]
      \item Roaring.io:s data är utdaterad (uppdateras kvartalsvis från Bolagsverket)
      \item Nyligen gjord ägarförändring (ännu ej registrerad)
      \item Indirekta ägare via holdingbolag (Roaring visar bara direktägande)
      \item Användaren har missförstått definitionen av "verklig huvudman" (PTL 1 kap 3-6 §§)
    \end{itemize}
  \item \textbf{Hur hanterar vi det?}
    \begin{itemize}[noitemsep]
      \item Flagga onboarding för manuell granskning (\texttt{requires\_manual\_review = true})
      \item Logga händelse med detaljer (personnr från form, VH från Roaring)
      \item Backend returnerar varning i response (men BLOCKERAR INTE processen)
      \item Admin Dashboard visar flaggade fall för manuell kontroll
      \item Redovisningskonsult kontaktar kunden för klargörande
    \end{itemize}
  \item \textbf{Omvänd situation:} Användaren svarar "Nej" men Roaring visar dem som VH >25\%
    \begin{itemize}[noitemsep]
      \item Kan vara legitimt (ägarandelen precis under 25\%, eller indirekt ägande)
      \item Flagga OCKSÅ för manuell granskning
      \item LOG event: "VH\_REVERSE\_MISMATCH\_DETECTED"
    \end{itemize}
\end{itemize}

\textbf{Trigger-definition och koppling till Celestial Risk Engine v3.0:}

\textbf{VIKTIGT: Fullständig metodbeskrivning finns i} \texttt{docs/Theory/metod\_riskbedömning\_kund\_v3.tex}

Tidigare version (Hot×Sårbarhet-modell) har ersatts med \textbf{Celestial Risk Engine 0–100 poäng-modell}:

\begin{enumerate}[noitemsep]
  \item \textbf{staticKycScore} (0–100): Beräknas från tre komponenter
    \begin{itemize}[noitemsep]
      \item Offentliga register (max 30 poäng): Bolagsålder, PEP, sanktioner, näringsförbud
      \item Korsvalidering (max 30 poäng): VH-mismatch, adressmismatch, F-skatt/moms-avvikelser
      \item Bokföringsdata (max 40 poäng): Sätts till 0 i Steg 1, fylls i vid SIE-analys
    \end{itemize}
  
  \item \textbf{sniHotFactor} (1–4): Hämtas från SNI-kod via Bolagsverket
    \begin{itemize}[noitemsep]
      \item 1 = Lågrisk (IT, konsult, design)
      \item 2 = Normal risk (handel, tillverkning)
      \item 3 = Förhöjd risk (restaurang, taxi, kontantintensiv)
      \item 4 = Hög risk (spelverksamhet, valutaväxling, ädelmetaller)
    \end{itemize}
  
  \item \textbf{adjustedScore}: staticKycScore + hotBonus
    \begin{itemize}[noitemsep]
      \item Hot 1: score + 0 poäng
      \item Hot 2: score + 10 poäng
      \item Hot 3: score + 20 poäng
      \item Hot 4: score + 30 poäng
    \end{itemize}
  
  \item \textbf{riskLevel} (fyrgradig skala enligt FFFS 2009:1):
    \begin{itemize}[noitemsep]
      \item \texttt{low}: adjustedScore 0–40
      \item \texttt{medium}: adjustedScore 41–70
      \item \texttt{high}: adjustedScore 71–90
      \item \texttt{critical}: adjustedScore > 90 ELLER hårda triggers
    \end{itemize}
  
  \item \textbf{Hårda triggers (automatisk critical-klassning):}
    \begin{itemize}[noitemsep]
      \item \texttt{hasSanctions = true} → STOPP, avslag av uppdrag
      \item \texttt{hasBusinessProhibition = true} → STOPP, avslag av uppdrag
      \item Kombination: VH-mismatch + saknad F-skatt + PEP → kritisk granskning
    \end{itemize}
  
  \item \textbf{Eskaleringstriggrar (höjer riskklass ett steg om på gräns):}
    \begin{itemize}[noitemsep]
      \item \texttt{isPEP = true} + score > 35 → Low blir Medium
      \item \texttt{vhMismatch = true} + score > 60 → Medium blir High
      \item \texttt{addressMismatch = true} + \texttt{highRiskIndustry = true} → Betoning
    \end{itemize}
\end{enumerate}

\textbf{Kostnad externa API-anrop:}
\begin{itemize}[noitemsep]
  \item \textbf{Best Case (2 VH direktägande):} Bolagsverket (1) + SPAR (1) + Roaring BO (1) + (2 VH × 3 checks) = 9 anrop
  \item \textbf{Worst Case (alternativ VH):} Bolagsverket (1) + SPAR (1) + Roaring BO (2) + (1 VH × 3 checks) = 7 anrop
\end{itemize}

% ------------------------------------------------------------------------
\subsection{Extern Endpoint 1: POST /organisationer (Bolagsverket)}
% ------------------------------------------------------------------------

\subsubsection{Beskrivning}

Validerar organisationsnummer och hämtar grundläggande företagsinformation inklusive VD och firmatecknare från Bolagsverkets officiella register.

\textbf{Fullständig URL:} \texttt{https://portal.api.bolagsverket.se/open/v1/companies}

\subsubsection{Request}

\textbf{Headers:}
\begin{lstlisting}
{
  "Content-Type": "application/json"
}
\end{lstlisting}

\textbf{Body (pseudokod):}
\begin{lstlisting}
{
  "organisationsnummer": ["556903-8671"]
}
\end{lstlisting}

\subsubsection{Response 200 OK}
\begin{lstlisting}
{
  "organisationer": [
    {
      "organisationsnummer": "string",
      "namn": "string",
      "organisationsform": "string (AB, HB, etc.)",
      "sniKod": "string",
      "sniBeskrivning": "string",
      "registreringsdatum": "ISO-8601 date",
      "status": "string (Aktiv, Avregistrerad)",
      "vd": {
        "namn": "string",
        "personnummer": "string"
      },
      "firmatecknare": [
        {
          "namn": "string",
          "personnummer": "string",
          "firmateckningsregel": "string"
        }
      ]
    }
  ]
}
\end{lstlisting}

\subsubsection{Implementation Notes för Response}

\textbf{Validering (pseudokod):}
\begin{lstlisting}
1. Extrahera organisationer array från response
2. OM tom array -> REJECT med 404 Not Found
3. Extrahera första org från array (endast en förväntad)
4. OM org.status !== "Aktiv" -> REJECT med 400 Bad Request
5. Extrahera VD personnummer: vd_personnr = org.vd.personnummer
6. Extrahera firmatecknare personnummer: firmatecknare_list = org.firmatecknare.map(ft => ft.personnummer)
7. Jämför användaren angett personnummer mot VD och firmatecknare:
   OM (user_personnr === vd_personnr OR user_personnr IN firmatecknare_list)
     -> GODKÄND
   ANNARS
     -> REJECT med 403 Forbidden ("Du är inte VD eller firmatecknare")
8. Spara hela org-objektet i onboarding_processes.bolagsverket_data
9. Spara hela response i external_api_responses tabell för audit trail
\end{lstlisting}

% ------------------------------------------------------------------------
\subsection{Extern Endpoint 2: POST /navet/spar/personnummer (Skatteverket SPAR)}
% ------------------------------------------------------------------------

\subsubsection{Beskrivning}

Validerar personnummer och hämtar folkbokföringsuppgifter från Skatteverkets SPAR-register (Svenska PersonAdressRegistret).

\textbf{Fullständig URL:} \texttt{https://navet.skatteverket.se/spar/personnummer}

\textbf{OBS:} Kräver X.509-certifikat för autentisering (test-certifikat för utveckling, prod-certifikat för produktion).

\subsubsection{Request}

\textbf{Headers:}
\begin{lstlisting}
{
  "Content-Type": "application/json"
}
\end{lstlisting}

\textbf{Body (pseudokod):}
\begin{lstlisting}
{
  "personnummer": "195001011234",
  "syfte": "Kundkännedom enligt PTL"
}
\end{lstlisting}

\subsubsection{Response 200 OK}
\begin{lstlisting}
{
  "personnummer": "string",
  "namn": {
    "fornamn": "string",
    "efternamn": "string"
  },
  "folkbokforingsadress": {
    "adress": "string",
    "postnummer": "string",
    "postort": "string"
  },
  "status": "string (Folkbokförd, Avregistrerad, Emigrerad)",
  "fodelsedatum": "ISO-8601 date"
}
\end{lstlisting}

\subsubsection{Implementation Notes för Response}

\textbf{Validering (pseudokod):}
\begin{lstlisting}
1. OM response.status !== "Folkbokförd" -> VARNING (inte blockerande, men flagga för manuell granskning)
2. Extrahera fullständigt namn: fullName = fornamn + " " + efternamn
3. Jämför namn mot Bolagsverket VD-namn (fuzzy match för att hantera skiljetecken):
   OM (fuzzyMatch(sparNamn, bolagsverketVdNamn) < 80% match)
     -> VARNING (namndiskrepans - flagga för manuell granskning)
4. Spara hela SPAR-objektet i onboarding_processes.spar_data
5. Spara hela response i external_api_responses tabell för audit trail
\end{lstlisting}

% ------------------------------------------------------------------------
\subsection{Extern Endpoint 3: GET /se/company/beneficial-owners/2.1/\{orgNr\} (Roaring)}
% ------------------------------------------------------------------------

\subsubsection{Beskrivning}

Hämtar registrerade verkliga huvudmän från Bolagsverkets register över verkliga huvudmän. Returnerar endast personer som äger/kontrollerar >25\% enligt PTL 3 kap 6 §.

\textbf{Fullständig URL:} \texttt{https://api.roaring.io/se/company/beneficial-owners/2.1/\{orgNr\}}

\textbf{PTL-grund:} PTL 3 kap 6 § - Identifiering av verklig huvudman

\subsubsection{Request}

\textbf{Headers:}
\begin{lstlisting}
{
  "Authorization": "Bearer <roaring_access_token>"
}
\end{lstlisting}

\subsubsection{Response 200 OK - Med VH >25\%}

\begin{lstlisting}
{
  "organizationNumber": "string",
  "companyName": "string",
  "beneficialOwners": [
    {
      "personalNumber": "string",
      "firstName": "string",
      "lastName": "string",
      "ownershipPercentage": float,
      "votingRightsPercentage": float,
      "ownershipType": "string (DIRECT, INDIRECT)",
      "source": "string",
      "verified": boolean
    }
  ],
  "totalCoverage": float
}
\end{lstlisting}

\subsubsection{Response 200 OK - Tom array (ingen >25\%)}
\begin{lstlisting}
{
  "organizationNumber": "string",
  "companyName": "string",
  "beneficialOwners": [],
  "totalCoverage": 0.0
}
\end{lstlisting}

\subsubsection{Implementation Notes för Response}

\textbf{Två-stegs algoritm (pseudokod):}
\begin{lstlisting}
1. Anropa Primary Endpoint: GET /beneficial-owners/2.1/{orgNr}
2. OM beneficialOwners.length > 0:
     -> Använd dessa VH (registrerade ägare >25%)
     -> Gå till steg 8
3. OM beneficialOwners.length === 0:
     -> Ingen ägare >25%, använd fallback (PTL 3 kap 8 paragraf)
     -> Anropa Alternative Endpoint: GET /alternative-beneficial-owner/1.0/{orgNr}
4. OM alternative response innehåller VD:
     -> Använd VD som alternativ VH
     -> Gå till steg 8
5. OM både primary OCH alternative ger tomt resultat:
     -> KRITISK ERROR (ingen VH identifierbar)
     -> REJECT onboarding
6. FÖR varje VH (från primary ELLER alternative):
     a. GET /pep/1.0/{personnr}
     b. GET /sanctions/1.0/{personnr}
     c. GET /businessprohibition/1.0/person/{personnr}
7. Analysera PEP/Sanctions/BusinessProhibition:
   OM isSanctioned = true -> REJECT omedelbart
   OM hasBusinessProhibition = true -> REJECT omedelbart
   OM isPEP = true -> Sätt trigger showPEP = true
8. Spara alla VH i onboarding_processes.beneficial_owners_data
9. Spara alla responses i external_api_responses för audit trail
\end{lstlisting}

% ------------------------------------------------------------------------
\subsection{Extern Endpoint 4: GET /se/company/alternative-beneficial-owner/1.0/\{orgNr\} (Roaring)}
% ------------------------------------------------------------------------

\subsubsection{Beskrivning}

Hämtar alternativ verklig huvudman (VD/högsta befattningshavare) när ingen ägare når 25\%-tröskeln enligt PTL 3 kap 8 §.

\textbf{Fullständig URL:} \texttt{https://api.roaring.io/se/company/alternative-beneficial-owner/1.0/\{orgNr\}}

\textbf{PTL-grund:} PTL 3 kap 8 § - Alternativ verklig huvudman (högsta befattningshavare)

\subsubsection{Request}

\textbf{Headers:}
\begin{lstlisting}
{
  "Authorization": "Bearer <roaring_access_token>"
}
\end{lstlisting}

\subsubsection{Response 200 OK}
\begin{lstlisting}
{
  "organizationNumber": "string",
  "companyName": "string",
  "alternativeBeneficialOwner": {
    "personalNumber": "string",
    "firstName": "string",
    "lastName": "string",
    "position": "string (VD, Styrelseordförande, etc.)",
    "source": "string",
    "verified": boolean
  }
}
\end{lstlisting}

\subsubsection{Implementation Notes för Response}

\textbf{Användning (pseudokod):}
\begin{lstlisting}
1. Anropas endast OM primary endpoint returnerar tom array
2. Extrahera alternativeBeneficialOwner från response
3. Verifiera att personnummer finns
4. Kör SAMMA PEP/Sanctions/BusinessProhibition checks som för vanlig VH
5. Markera i databasen att denna VH är "alternativ" (för rapportering)
\end{lstlisting}

% ------------------------------------------------------------------------
\subsection{Extern Endpoint 5: GET /se/pep/1.0/\{personnr\} (Roaring)}
% ------------------------------------------------------------------------

\subsubsection{Beskrivning}

Kontrollerar om person är Politically Exposed Person (PEP) enligt PTL 3 kap 10 §. Inkluderar nationella och internationella PEP-listor samt familjemedlemmar/nära medarbetare.

\textbf{Fullständig URL:} \texttt{https://api.roaring.io/se/pep/1.0/\{personnr\}}

\subsubsection{Request}

\textbf{Headers:}
\begin{lstlisting}
{
  "Authorization": "Bearer <roaring_access_token>"
}
\end{lstlisting}

\subsubsection{Response 200 OK - Är PEP}

\begin{lstlisting}
{
  "personalNumber": "string",
  "isPEP": true,
  "pepType": "string (DOMESTIC, FOREIGN, RCA)",
  "position": "string",
  "organization": "string",
  "activeSince": "ISO-8601 date",
  "activeUntil": "ISO-8601 date (nullable)"
}
\end{lstlisting}

\subsubsection{Response 200 OK - Inte PEP}
\begin{lstlisting}
{
  "personalNumber": "string",
  "isPEP": false
}
\end{lstlisting}

\subsubsection{Implementation Notes för Response}

\textbf{Hantering (pseudokod):}
\begin{lstlisting}
1. OM isPEP = true:
     → Sätt trigger showPEP = true
     → Spara PEP-detaljer för fördjupade frågor i Steg 2
     → Fortsätt onboarding (PEP är INTE blockerande, kräver bara EDD)
2. OM isPEP = false:
     → Ingen ytterligare åtgärd
\end{lstlisting}

% ------------------------------------------------------------------------
\subsection{Extern Endpoint 6: GET /se/sanctions/1.0/\{personnr\} (Roaring)}
% ------------------------------------------------------------------------

\subsubsection{Beskrivning}

Kontrollerar om person finns på sanktionslistor (EU, UN, OFAC, etc.). Detta är BLOCKERANDE - om personen är sanktionerad MÅSTE onboarding avbrytas omedelbart.

\textbf{Fullständig URL:} \texttt{https://api.roaring.io/se/sanctions/1.0/\{personnr\}}

\subsubsection{Request}

\textbf{Headers:}
\begin{lstlisting}
{
  "Authorization": "Bearer <roaring_access_token>"
}
\end{lstlisting}

\subsubsection{Response 200 OK - Är sanktionerad}

\begin{lstlisting}
{
  "personalNumber": "string",
  "isSanctioned": true,
  "sanctionLists": [
    {
      "listName": "string (EU Sanctions, UN Sanctions, etc.)",
      "addedDate": "ISO-8601 date",
      "reason": "string"
    }
  ]
}
\end{lstlisting}

\subsubsection{Response 200 OK - Inte sanktionerad}
\begin{lstlisting}
{
  "personalNumber": "string",
  "isSanctioned": false
}
\end{lstlisting}

\subsubsection{Implementation Notes för Response}

\textbf{Hantering (pseudokod):}
\begin{lstlisting}
1. OM isSanctioned = true:
     → REJECT onboarding OMEDELBART
     → Returnera 403 Forbidden med message: "Sanktionerad person - onboarding ej möjlig"
     → Logga incident för compliance-team
     → AVSLUTA
2. OM isSanctioned = false:
     → Fortsätt onboarding
\end{lstlisting}

% ------------------------------------------------------------------------
\subsection{Extern Endpoint 7: GET /se/businessprohibition/1.0/person/\{personnr\} (Roaring)}
% ------------------------------------------------------------------------

\subsubsection{Beskrivning}

Kontrollerar om person har näringsförbud enligt bokföringslagen. Näringsförbud innebär att personen inte får driva företag och är BLOCKERANDE för onboarding.

\textbf{Fullständig URL:} \texttt{https://api.roaring.io/se/businessprohibition/1.0/person/\{personnr\}}

\subsubsection{Request}

\textbf{Headers:}
\begin{lstlisting}
{
  "Authorization": "Bearer <roaring_access_token>"
}
\end{lstlisting}

\subsubsection{Response 200 OK - Har näringsförbud}

\begin{lstlisting}
{
  "personalNumber": "string",
  "hasBusinessProhibition": true,
  "prohibitionStart": "ISO-8601 date",
  "prohibitionEnd": "ISO-8601 date",
  "reason": "string"
}
\end{lstlisting}

\subsubsection{Response 200 OK - Inget näringsförbud}
\begin{lstlisting}
{
  "personalNumber": "string",
  "hasBusinessProhibition": false
}
\end{lstlisting}

\subsubsection{Implementation Notes för Response}

\textbf{Hantering (pseudokod):}
\begin{lstlisting}
1. OM hasBusinessProhibition = true:
     → REJECT onboarding OMEDELBART
     → Returnera 403 Forbidden med message: "Näringsförbud - onboarding ej möjlig"
     → Logga incident för compliance-team
     → AVSLUTA
2. OM hasBusinessProhibition = false:
     → Fortsätt onboarding
\end{lstlisting}

% ------------------------------------------------------------------------
\subsection{Backend Persistent Lagring}
% ------------------------------------------------------------------------

\textbf{Tabell: onboarding\_processes}
\begin{lstlisting}[language=sql]
-- Utöka befintlig tabell med kolumner för Steg 1
ALTER TABLE onboarding_processes
ADD COLUMN IF NOT EXISTS riskfragor_steg1 JSONB,
ADD COLUMN IF NOT EXISTS bolagsverket_data JSONB,
ADD COLUMN IF NOT EXISTS spar_data JSONB,
ADD COLUMN IF NOT EXISTS beneficial_owners_data JSONB,
ADD COLUMN IF NOT EXISTS triggers JSONB DEFAULT '{}'::jsonb,
ADD COLUMN IF NOT EXISTS red_flags JSONB DEFAULT '[]'::jsonb,
ADD COLUMN IF NOT EXISTS vh_mismatch_detected BOOLEAN DEFAULT false,
ADD COLUMN IF NOT EXISTS vh_mismatch_reason TEXT,
ADD COLUMN IF NOT EXISTS requires_manual_review BOOLEAN DEFAULT false;

-- Index för snabba sökningar
CREATE INDEX IF NOT EXISTS idx_onboarding_triggers 
  ON onboarding_processes USING GIN (triggers);
CREATE INDEX IF NOT EXISTS idx_onboarding_red_flags 
  ON onboarding_processes USING GIN (red_flags);
CREATE INDEX IF NOT EXISTS idx_manual_review
  ON onboarding_processes (requires_manual_review)
  WHERE requires_manual_review = true;
\end{lstlisting}

\textbf{Tabell: external\_api\_responses}
\begin{lstlisting}[language=sql]
-- Logga alla externa API-anrop för audit trail och GDPR
CREATE TABLE IF NOT EXISTS external_api_responses (
  id UUID PRIMARY KEY DEFAULT gen_random_uuid(),
  onboarding_id UUID NOT NULL REFERENCES onboarding_processes(id),
  api_name VARCHAR(100) NOT NULL,  -- "bolagsverket", "spar", "roaring_pep", etc.
  endpoint VARCHAR(255) NOT NULL,   -- "/organisationer", "/pep/1.0/{personnr}", etc.
  request_data JSONB,               -- Request body (om relevant)
  response_data JSONB,              -- Fullständig response
  response_code INTEGER,            -- HTTP status code
  created_at TIMESTAMP DEFAULT NOW()
);

-- Index för snabba sökningar
CREATE INDEX IF NOT EXISTS idx_api_responses_onboarding 
  ON external_api_responses(onboarding_id);
CREATE INDEX IF NOT EXISTS idx_api_responses_api_name 
  ON external_api_responses(api_name);
CREATE INDEX IF NOT EXISTS idx_api_responses_created 
  ON external_api_responses(created_at);
\end{lstlisting}

\textbf{Update-query (pseudokod):}
\begin{lstlisting}[language=sql]
UPDATE onboarding_processes
SET 
  riskfragor_steg1 = $request_body,
  bolagsverket_data = $bolagsverket_response,
  spar_data = $spar_response,
  beneficial_owners_data = $roaring_beneficial_owners,
  triggers = $calculated_triggers,
  red_flags = $detected_red_flags,
  current_step = 2,
  status = 'in_progress',
  updated_at = NOW()
WHERE id = $onboardingId
  AND user_id = $user_id_from_jwt;
\end{lstlisting}

% ------------------------------------------------------------------------
\subsection{Frontend Datalagring i Browser}
% ------------------------------------------------------------------------

\subsubsection{localStorage}

\begin{itemize}[noitemsep]
  \item \texttt{onboardingId} (UUID) — Identifierar processen, sätts av POST /uppdrag
  \item \texttt{currentStep} (Integer) — Vilket steg användaren är på (0-7)
  \item \texttt{riskfragorSteg1Draft} (JSON) — Autosave medan användaren fyller i formulär
  \item \texttt{triggers} (JSON) — Vilka riskflaggor som är aktiva (används för UI-betoning)
  \item \texttt{wizardConfig} (JSON) — Wizard-steg från backend (default + custom från config.json)
  \item \texttt{bolagsverketData} (JSON) — Cached företagsinformation för visning
  \item \texttt{sparData} (JSON) — Cached personuppgifter för visning
\end{itemize}

\textbf{Exempel localStorage efter Steg 1 (utan custom config):}
\begin{lstlisting}[language=json]
{
  "onboardingId": "550e8400-e29b-41d4-a716-446655440000",
  "currentStep": 2,
  "triggers": {
    "showUtlandska": true,
    "showKontanter": false,
    "showVerksamhet": false,
    "showPEP": false
  },
  "wizardConfig": {
    "totalSteps": 4,
    "hasCustomConfig": false,
    "steps": [
      {"id": 2, "title": "Utländska transaktioner", "isCustom": false},
      {"id": 3, "title": "Kunder & Affärspartners", "isCustom": false},
      {"id": 4, "title": "Betalningar & Transaktioner", "isCustom": false}
    ]
  },
  "bolagsverketData": {
    "foretagsnamn": "Exempel AB",
    "sniKod": "47111",
    "sniBeskrivning": "Livsmedelsbutiker"
  }
}
\end{lstlisting}

\textbf{Exempel localStorage efter Steg 1 (med custom config från Inställningar):}
\begin{lstlisting}[language=json]
{
  "onboardingId": "550e8400-...",
  "currentStep": 2,
  "wizardConfig": {
    "totalSteps": 6,
    "hasCustomConfig": true,
    "customConfigVersion": "v2.1",
    "steps": [
      {"id": 2, "title": "Utländska transaktioner", "isCustom": false},
      {"id": 3, "title": "Kunder & Affärspartners", "isCustom": false},
      {"id": 4, "title": "Betalningar & Transaktioner", "isCustom": false},
      {"id": 5, "title": "Miljöhänsyn & Hållbarhet", "isCustom": true, "questionCount": 4},
      {"id": 6, "title": "IT-säkerhet & GDPR", "isCustom": true, "questionCount": 6}
    ]
  }
}
\end{lstlisting}

\subsubsection{Cookies}

\begin{itemize}[noitemsep]
  \item \texttt{access\_token} (HttpOnly, Secure, SameSite=Strict) — JWT access token för API-autentisering
  \item \texttt{refresh\_token} (HttpOnly, Secure, SameSite=Strict) — JWT refresh token för token renewal
\end{itemize}

Inga cookies specifika för Riskfrågor Steg 1.

% ------------------------------------------------------------------------
\subsection{Backend - Behandlingshistorik och GDPR}
% ------------------------------------------------------------------------

\subsubsection{Loggade händelser}

Alla API-anrop och viktiga händelser loggas för GDPR-compliance (rätt till tillgång) och audit trail.

\begin{itemize}[noitemsep]
  \item \texttt{riskfragor\_steg1\_started} — När användaren besöker sliden
  \item \texttt{riskfragor\_steg1\_autosave} — När formulär auto-sparas (localStorage)
  \item \texttt{riskfragor\_steg1\_submitted} — När användaren klickar \enquote{Nästa}
  \item \texttt{bolagsverket\_api\_called} — Organisationsnummer validering
  \item \texttt{spar\_api\_called} — Personnummer validering
  \item \texttt{roaring\_beneficial\_owner\_primary\_called} — VH >25\% check
  \item \texttt{roaring\_beneficial\_owner\_alternative\_called} — Alternativ VH check (om behövs)
  \item \texttt{vh\_mismatch\_detected} — \textbf{(NYT)} Användaren hävdar VH men finns ej i Roaring
  \item \texttt{vh\_reverse\_mismatch\_detected} — \textbf{(NYT)} Roaring visar VH men användaren nekar
  \item \texttt{manual\_review\_flagged} — \textbf{(NYT)} Ärendet flaggat för manuell granskning
  \item \texttt{roaring\_pep\_called} — PEP-kontroll per VH
  \item \texttt{roaring\_sanctions\_called} — Sanktionskontroll per VH
  \item \texttt{roaring\_businessprohibition\_called} — Näringsförbud per VH
  \item \texttt{triggers\_calculated} — Vilka fördjupningsfrågor triggades
  \item \texttt{red\_flags\_detected} — Om PEP/Sanctions/BusinessProhibition hittades
  \item \texttt{onboarding\_rejected} — Om sanktionerad eller näringsförbud
\end{itemize}

\subsubsection{CSV-filer (Primär databasimplementation)}

\textbf{UTVECKLINGSSTRATEGI:} Hela applikationen byggs FÖRST med CSV-filer för ALL databashantering. PostgreSQL införs som slutoptimering ENDAST efter att hela applikationen fungerar felfritt. Detta undviker race conditions och concurrent write-problem under utveckling.

\textbf{Varför CSV först?}
\begin{itemize}[noitemsep]
  \item Enklare debugging (kan öppna CSV i editor och se exakt vad som sparats)
  \item Inga databaskonfigurationsproblem under utveckling
  \item Snabbare iterationer (ingen migration/schema-hantering)
  \item Fungerar perfekt för single-user development environment
  \item PostgreSQL blir "drop-in replacement" när applikationen är stabil
\end{itemize}

\textbf{onboarding\_processes.csv:}
\begin{lstlisting}
id,user_id,riskfragor_steg1,triggers,vh_mismatch_detected,requires_manual_review,current_step,created_at,updated_at
550e8400-...,123e4567-...,"{\"affarsIde\":\"...\"}","{\"showUtlandska\":true}",false,false,2,2025-10-24T10:00:00Z,2025-10-24T10:05:00Z
\end{lstlisting}

\textbf{external\_api\_responses.csv:}
\begin{lstlisting}
id,onboarding_id,api_name,endpoint,response_code,response_body,created_at
uuid1,550e8400-...,bolagsverket,/organisationer,200,"{\"foretagsnamn\":\"...\"}",2025-10-24T10:01:00Z
uuid2,550e8400-...,spar,/navet/spar/personnummer,200,"{\"namn\":\"...\"}",2025-10-24T10:02:00Z
uuid3,550e8400-...,roaring_pep,/pep/1.0/195001011234,200,"{\"isPEP\":false}",2025-10-24T10:03:00Z
\end{lstlisting}

\textbf{event\_log.csv:}
\begin{lstlisting}
id,onboarding_id,event_type,user_id,details,created_at
uuid1,550e8400-...,riskfragor_steg1_submitted,123e4567-...,"{\"step\":1}",2025-10-24T10:05:00Z
uuid2,550e8400-...,vh_mismatch_detected,123e4567-...,"{\"reason\":\"...\"}",2025-10-24T10:05:30Z
\end{lstlisting}

% ------------------------------------------------------------------------
\subsection{Databastabeller (Framtida PostgreSQL-optimering)}
% ------------------------------------------------------------------------

\textbf{OBS:} SQL-scheman nedan dokumenterar framtida PostgreSQL-implementation. NUVARANDE implementation använder CSV-filer med samma struktur. PostgreSQL införs som slutoptimering för multi-user concurrent access.
Se sektion "Backend Persistent Lagring" ovan för fullständigt SQL-schema.

\textbf{Sammanfattning:}
\begin{itemize}[noitemsep]
  \item \texttt{onboarding\_processes} — Huvudtabell, innehåller riskfragor\_steg1, triggers, bolagsverket\_data, spar\_data, beneficial\_owners\_data, red\_flags
  \item \texttt{external\_api\_responses} — Audit trail för alla externa API-anrop, sparar request + response för GDPR-compliance
  \item \texttt{user\_settings} — Innehåller custom\_wizard\_config (JSON) per användare
  \item Index på triggers (GIN för JSONB) och red\_flags (GIN för JSONB) för snabba sökningar
\end{itemize}

% ------------------------------------------------------------------------
\subsection{Custom Wizard Config (Valfri anpassning via Inställningar)}
% ------------------------------------------------------------------------

\subsubsection{Översikt}

Användare kan ladda upp en custom \texttt{config.json} via Inställningar-sidan för att lägga till EGNA wizard-steg (Steg 5, 6, 7...) med anpassade frågor. De obligatoriska Steg 2-4 kan INTE ersättas eller tas bort, endast kompletteras.

\textbf{När läses config.json?} Backend läser användarens custom config vid \texttt{POST /riskfragor/steg1} (precis innan Steg 2 ska visas). Detta är den ENKLASTE approachen — ingen extra API-anrop, ingen cache-invalidation, alltid uppdaterad config.

\subsubsection{Upload Endpoint}

\textbf{Endpoint:} \texttt{POST /api/user/settings/config}

\textbf{Request:}
\begin{lstlisting}
{
  "configFile": "base64-encoded JSON eller multipart/form-data",
  "version": "v2.1"  // valfri versionstagg
}
\end{lstlisting}

\textbf{Response 200 OK:}
\begin{lstlisting}
{
  "success": true,
  "message": "Custom wizard config sparad",
  "totalStepsAdded": 2,
  "configVersion": "v2.1"
}
\end{lstlisting}

\textbf{Backend sparar:}
\begin{lstlisting}[language=sql]
-- CSV: user_settings.csv
-- SQL (framtida):
UPDATE user_settings
SET custom_wizard_config = $1::jsonb,
    config_version = $2,
    updated_at = NOW()
WHERE user_id = $3;
\end{lstlisting}

\subsubsection{Config.json Structure}

\begin{lstlisting}[language=json]
{
  "version": "v2.1",
  "additionalSteps": [
    {
      "id": 5,
      "title": "Miljöhänsyn & Hållbarhet",
      "description": "Frågor om företagets miljöarbete",
      "conditional": false,
      "questions": [
        {
          "id": "env_1",
          "text": "Har företaget miljöcertifiering (t.ex. ISO 14001)?",
          "type": "radio",
          "required": true,
          "options": ["Ja", "Nej", "Under arbete"]
        },
        {
          "id": "env_2",
          "text": "Beskriv företagets miljöpolicy (om tillämpligt)",
          "type": "textarea",
          "required": false,
          "maxLength": 500,
          "placeholder": "T.ex. återvinning, koldioxidkompensation..."
        },
        {
          "id": "env_3",
          "text": "Hållbarhetsrapportering enligt GRI/CSRD?",
          "type": "checkbox",
          "required": false,
          "label": "Ja, vi rapporterar enligt standarder"
        }
      ]
    },
    {
      "id": 6,
      "title": "IT-säkerhet & GDPR",
      "description": "Frågor om informationssäkerhet",
      "conditional": false,
      "questions": [
        {
          "id": "it_1",
          "text": "Har företaget utsett ett dataskyddsombud (DPO)?",
          "type": "radio",
          "required": true,
          "options": ["Ja", "Nej", "Ej tillämpligt (< 250 anställda)"]
        },
        {
          "id": "it_2",
          "text": "Genomförs regelbundna säkerhetsrevisioner?",
          "type": "radio",
          "required": true,
          "options": ["Ja, årligen", "Ja, vartannat år", "Nej", "Planerat"]
        }
      ]
    }
  ]
}
\end{lstlisting}

\subsubsection{Regler för Custom Config}

\begin{itemize}[noitemsep]
  \item \textbf{Steg 2-4 är obligatoriska} — Custom steg får ENDAST id 5 och uppåt
  \item \textbf{Max 10 additional steps} — För att undvika orimligt långa wizards
  \item \textbf{Max 20 questions per step} — För att hålla steg hanterliga
  \item \textbf{Conditional steg} — Kan använda triggers från Steg 1 (men rekommenderas ej för enkelhet)
  \item \textbf{Question types} — Stödda typer: radio, checkbox, textarea, dropdown, number
  \item \textbf{Validation} — Backend validerar config.json schema vid upload
  \item \textbf{Fallback} — Om custom config är invalid/korrupt, använd default Steg 2-4 bara
\end{itemize}

\subsubsection{Frontend Rendering}

Frontend använder \texttt{wizardConfig} från POST /riskfragor/steg1 response för att dynamiskt rendera steg:

\begin{lstlisting}[language=javascript]
// RiskFragorWizard.jsx
const { wizardConfig } = responseData;

{wizardConfig.steps.map((step, index) => (
  <WizardStep 
    key={step.id}
    stepNumber={step.id}
    totalSteps={wizardConfig.totalSteps}
    title={step.title}
    questions={step.questions || DEFAULT_QUESTIONS[step.id]}
    isCustom={step.isCustom}
  />
))}
\end{lstlisting}

% ------------------------------------------------------------------------
\subsection{Sammanfattning (Prosaisk beskrivning)}
% ------------------------------------------------------------------------

\subsubsection{Scenario 1: Första gången - Inga riskflaggor}

\textbf{Aktörer:}
\begin{itemize}[noitemsep]
  \item \textbf{Användare (inloggad):} Anna Andersson, redovisningskonsult på byrån (hennes userId i JWT token)
  \item \textbf{Presumtiv klient:} Celestial AB (org.nr 556903-8671), VD Bengt Bengtsson (personnr 196001011234)
\end{itemize}

Anna Andersson (redovisningskonsulten) sitter och intervjuar Bengt Bengtsson, VD för Celestial AB som vill bli kund. Anna har öppnat onboarding-appen och kommer från Uppdragsval-sliden med \texttt{onboardingId} sparat i localStorage. 

På Riskfrågor Steg 1-sliden välkomnas hon med en intro-text om varför riskfrågor krävs (PTL-referens). Hon fyller i 9 grundfrågor baserat på information från Bengt:

\begin{enumerate}[noitemsep]
  \item \textbf{Affärsidé:} "Restaurang med mat från medelhavet"
  \item \textbf{Kundtyper:} Privatpersoner OCH företag (både walk-in kunder och cateringuppdrag)
  \item \textbf{Utländska partners:} "Nej"
  \item \textbf{Leverantörer:} "Menigo, Martin \& Servera, Systembolaget"
  \item \textbf{Verksamhet ändrad:} "Nej"
  \item \textbf{Organisationsnummer:} 556903-8671 (Celestial AB)
  \item \textbf{Personnummer VD:} 196001011234 (Bengt Bengtsson)
  \item \textbf{Verklig huvudman:} "Nej" (Bengt äger bara 15\%, inte >25\%)
  \item \textbf{PEP-status:} Avmarkerad
\end{enumerate}

När Anna klickar \enquote{Nästa} visas en loading spinner med text \enquote{Validerar uppgifter mot externa register...}. Backend gör följande:

\textbf{Steg 1-3: Autentisering och validering}
\begin{itemize}[noitemsep]
  \item Extraherar user\_id från JWT token (Anna Andersson's userId - redovisningskonsulten)
  \item Verifierar att onboardingId tillhör Anna (security check)
  \item Validerar format på org.nr (556903-8671) och personnr (196001011234)
\end{itemize}

\textbf{Steg 4: Bolagsverket API-anrop}

Backend anropar Bolagsverket POST /organisationer med org.nr 556903-8671 (klientens företag). Response innehåller:
\begin{lstlisting}
{
  "organisationsnummer": "556903-8671",
  "namn": "Celestial AB",
  "organisationsform": "AB",
  "sniKod": "56101",
  "sniBeskrivning": "Restauranger och gatukök",
  "vd": {
    "namn": "Bengt Bengtsson",
    "personnummer": "196001011234"
  },
  "firmatecknare": [
    {
      "namn": "Bengt Bengtsson",
      "personnummer": "196001011234"
    }
  ]
}
\end{lstlisting}

Backend verifierar att Bengt's personnummer (196001011234) matchar VD-personnumret från Bolagsverket. GODKÄND

\textbf{Steg 5: SPAR API-anrop}

Backend anropar Skatteverket SPAR POST /navet/spar/personnummer med personnummer 196001011234 (Bengt Bengtsson). Response innehåller:
\begin{lstlisting}
{
  "personnummer": "196001011234",
  "namn": {
    "fornamn": "Bengt",
    "efternamn": "Bengtsson"
  },
  "folkbokforingsadress": {
    "adress": "Storgatan 1",
    "postnummer": "12345",
    "postort": "Stockholm"
  },
  "status": "Folkbokförd"
}
\end{lstlisting}

Backend verifierar att Bengt är folkbokförd. GODKÄND

\textbf{Steg 6-7: Roaring Beneficial Owners (två-stegs algoritm)}

Backend anropar först Roaring Primary Endpoint GET /beneficial-owners/2.1/556903-8671 (Celestial AB). Response innehåller:
\begin{lstlisting}
{
  "organizationNumber": "5569038671",
  "companyName": "Celestial AB",
  "beneficialOwners": [
    {
      "personalNumber": "197001011234",
      "firstName": "Bengt",
      "lastName": "Bengtsson",
      "ownershipPercentage": 60.0,
      "votingRightsPercentage": 60.0
    },
    {
      "personalNumber": "196501011234",
      "firstName": "Carl",
      "lastName": "Carlsson",
      "ownershipPercentage": 40.0,
      "votingRightsPercentage": 40.0
    }
  ]
}
\end{lstlisting}

Två verkliga huvudmän identifierade (båda >25\%). Backend behöver INTE anropa Alternative Endpoint eftersom primary gav resultat.

\textbf{Steg 8: PEP/Sanctions/BusinessProhibition checks}

För varje VH (Bengt och Carl) anropar backend:
\begin{itemize}[noitemsep]
  \item GET /pep/1.0/197001011234 → isPEP: false
  \item GET /sanctions/1.0/197001011234 → isSanctioned: false
  \item GET /businessprohibition/1.0/person/197001011234 → hasBusinessProhibition: false
  \item GET /pep/1.0/196501011234 → isPEP: false
  \item GET /sanctions/1.0/196501011234 → isSanctioned: false
  \item GET /businessprohibition/1.0/person/196501011234 → hasBusinessProhibition: false
\end{itemize}

Alla checks ger grönt ljus. INGA RED FLAGS

\textbf{Steg 9: Beräkna riskAssessment (Celestial Risk Engine v3.0)}

\textbf{Komponent 1: Offentliga register (publicData)}
\begin{itemize}[noitemsep]
  \item Bolagsålder: Celestial AB registrerat 2023-01-15 = 2 år → +4 poäng
  \item PEP-status: Ingen VH är PEP → 0 poäng
  \item Sanktioner: Ingen VH på sanktionslista → 0 poäng
  \item Näringsförbud: Ingen VH har näringsförbud → 0 poäng
  \item \textbf{Totalt publicData:} 4 poäng (av max 30)
\end{itemize}

\textbf{Komponent 2: Korsvalidering (crossValidation)}
\begin{itemize}[noitemsep]
  \item VH-mismatch: Klienten svarade \enquote{Nej} (ej VH) men Roaring visar Bengt som 60\% ägare
    → Omvänd mismatch detekterad → +8 poäng
  \item Adressmismatch: SPAR visar Storgatan 1, Bolagsverket visar Storgatan 1 → 0 poäng
  \item F-skatt: (Kontrolleras ej i detta scenario) → 0 poäng
  \item Momsregistrering: (Kontrolleras ej i detta scenario) → 0 poäng
  \item \textbf{Totalt crossValidation:} 8 poäng (av max 30)
\end{itemize}

\textbf{Komponent 3: Bokföringsdata (bookkeeping)}
\begin{itemize}[noitemsep]
  \item Ännu ej analyserad i Steg 1 → 0 poäng (av max 40)
\end{itemize}

\textbf{Beräkning:}
\begin{align*}
\text{staticKycScore} &= 4 + 8 + 0 = 12 \text{ poäng} \\
\text{sniHotFactor} &= 3 \text{ (SNI 56101 = Restaurang, kontantintensiv)} \\
\text{hotBonus} &= (3-1) \times 10 = 20 \text{ poäng} \\
\text{adjustedScore} &= 12 + 20 = 32 \text{ poäng}
\end{align*}

\textbf{Riskkla ssificering:}
\begin{itemize}[noitemsep]
  \item adjustedScore = 32 → Inom intervall 0–40 → \texttt{riskLevel = "low"}
  \item Inga hårda triggers (sanctions/businessProhibition) → \texttt{requiresEDD = false}
  \item VH-mismatch detekterad → \texttt{requiresManualReview = true}
\end{itemize}

\textbf{Steg 10: Beräkna triggers för UI-betoning}

Backend sätter triggers baserat på klientens svar OCH API-resultat:
\begin{itemize}[noitemsep]
  \item isPEP = false (varken VD eller VH är PEP)
  \item hasSanctions = false
  \item hasBusinessProhibition = false
  \item vhMismatch = true (omvänd: klient säger nej men Roaring säger ja)
  \item addressMismatch = false
  \item missingFTax = false
  \item missingVAT = false
  \item newCompany = false (2 år gammalt)
  \item highRiskIndustry = true (SNI 56101 = restaurang)
\end{itemize}

\textbf{OBS:} Triggers används INTE för att hoppa över Steg 2-4. Steg 2-4 visas ALLTID (eller fler steg enligt config.json). Triggers används för att flagga risker och kan påverka vilka frågor som betonas, men wizard-stegen körs alltid igenom.

\textbf{Steg 10-11: Spara i databas och returnera}

Backend uppdaterar onboarding\_processes tabell:
\begin{lstlisting}[language=sql]
UPDATE onboarding_processes
SET 
  riskfragor_steg1 = '{...Bengts/Celestials svar...}',
  bolagsverket_data = '{...Celestial AB data...}',
  spar_data = '{...Bengts folkbokföringsdata...}',
  beneficial_owners_data = '[{...Bengt...}, {...Carl...}]',
  triggers = '{"showUtlandska":false,"showKontanter":false,"showVerksamhet":false,"showPEP":false}',
  current_step = 2,
  updated_at = NOW()
WHERE id = '550e8400-...' AND user_id = 'Anna_userId';
\end{lstlisting}

Backend sparar också alla 9 externa API-anrop i external\_api\_responses för audit trail.

Backend returnerar response:
\begin{lstlisting}
{
  "success": true,
  "triggers": {"showUtlandska":false,"showKontanter":true,"showVerksamhet":false,"showPEP":false},
  "message": "Grundläggande riskbedömning klar - kontanthantering flaggad",
  "nextStep": "/riskfragor/steg2"
}
\end{lstlisting}

\textbf{Frontend-hantering:}

Frontend tar emot response, sparar triggers i localStorage, och navigerar ALLTID till \texttt{/riskfragor/steg2}. Frontend läser sedan config.json för att se vilka wizard-steg som ska visas (minst Steg 2-4, eventuellt fler enligt användarens config.json). Triggers används för att flagga risker i backend och kan påverka vilka frågor som betonas i UI, men påverkar INTE om stegen visas eller inte.

\textbf{OBS om databas:} Under utveckling används CSV-filer för ALL databashantering (onboarding\_processes.csv, external\_api\_responses.csv, etc.). PostgreSQL SQL införs FÖRST när hela applikationen fungerar felfritt som slutoptimering för att hantera fler samtidiga användare utan skrivfel/race conditions.

\subsubsection{Scenario 2: Triggers aktiverade (högre risk)}

\textbf{Aktörer:}
\begin{itemize}[noitemsep]
  \item \textbf{Användare (inloggad):} Anna Andersson, redovisningskonsult på byrån
  \item \textbf{Presumtiv klient:} TechStart AB, VD Eva Eriksson
\end{itemize}

Anna intervjuar Eva Eriksson, VD för TechStart AB som vill bli kund. Eva's företag har annorlunda profil än Celestial AB:
\begin{itemize}[noitemsep]
  \item \textbf{Kundtyper:} Privatpersoner OCH företag (högre risk)
  \item \textbf{Utländska partners:} "Ja, vi har kunder i Norge och Danmark" (gränsöverskridande verksamhet)
\end{itemize}

När backend beräknar triggers:
\begin{itemize}[noitemsep]
  \item showUtlandska = true (utlandskaPartners !== "Nej")
  \item showKontanter = true (kundTyper.privatpersoner = true)
\end{itemize}

Backend returnerar:
\begin{lstlisting}
{
  "success": true,
  "triggers": {"showUtlandska":true,"showKontanter":true,"showVerksamhet":false,"showPEP":false},
  "message": "Riskflaggor identifierade (utländska transaktioner + kontanter)",
  "nextStep": "/riskfragor/steg2"
}
\end{lstlisting}

\textbf{Frontend-hantering:}

Precis som i Scenario 1, navigerar frontend ALLTID till \texttt{/riskfragor/steg2}. Skillnaden är att triggers nu är aktiva (\texttt{showUtlandska = true}, \texttt{showKontanter = true}), vilket kan användas för att:
\begin{itemize}[noitemsep]
  \item Betona specifika frågor i UI (t.ex. markera utländska transaktioner med varning)
  \item Visa extra hjälptext eller förklaringar för riskfyllda områden
  \item Flagga ärendet för manuell granskning i backend
  \item Påverka risk-score i slutlig bedömning
\end{itemize}

Frontend visar SAMMA wizard-steg (Steg 2-4) som i Scenario 1, men med triggers kan vissa frågor betonas eller flaggas för extra uppmärksamhet.

\textbf{Steg 2-4 visas ALLTID (konfigurerbart via config.json):}
\begin{enumerate}[noitemsep]
  \item \textbf{Steg 2 (Utländska transaktioner):} Vilka länder, omsättningsandel, fast driftställe
  \item \textbf{Steg 3 (Kunder \& Partners):} Återkommande partners, största leverantörer/kunder
  \item \textbf{Steg 4 (Betalningar):} Betalningsmetoder, kontantandel, stora transaktioner
\end{enumerate}

Efter att Anna besvarat Steg 2-4 (+ eventuella extra steg från config.json), sparas alla svar via \texttt{POST /riskfragor/steg2} (en enda batch), och frontend navigerar till \texttt{/identitetskontroll}.

\subsubsection{Scenario 3: Validering misslyckas - 403 Forbidden}

\textbf{Aktörer:}
\begin{itemize}[noitemsep]
  \item \textbf{Användare (inloggad):} Anna Andersson, redovisningskonsult på byrån
  \item \textbf{Presumtiv klient:} Celestial AB (org.nr 556903-8671), VD Bengt Bengtsson
  \item \textbf{Person som försöker representera företaget:} David Davidsson (personnr 198001011234) - INTE behörig
\end{itemize}

Anna intervjuar David Davidsson som påstår sig vara VD för Celestial AB. Hon fyller i formuläret:
\begin{itemize}[noitemsep]
  \item \textbf{Organisationsnummer:} 556903-8671 (Celestial AB)
  \item \textbf{Personnummer VD:} 198001011234 (David Davidsson)
\end{itemize}

Backend anropar Bolagsverket och får response som visar att Bengt Bengtsson (INTE David) är VD:
\begin{lstlisting}
{
  "vd": {
    "namn": "Bengt Bengtsson",
    "personnummer": "196001011234"
  },
  "firmatecknare": [
    {
      "namn": "Bengt Bengtsson",
      "personnummer": "196001011234"
    }
  ]
}
\end{lstlisting}

Backend jämför det angivna personnumret (198001011234 - David) med VD och firmatecknare från Bolagsverket. INGEN MATCH hittas - David är inte behörig att representera företaget.

Backend returnerar:
\begin{lstlisting}
{
  "success": false,
  "error": "not_authorized",
  "code": "PERSON_NOT_AUTHORIZED",
  "details": "Personnummer stämmer inte med VD/firmatecknare enligt Bolagsverket"
}
\end{lstlisting}

Frontend visar error message i röd banner: "Du är inte behörig att skapa uppdrag för detta bolag. Endast VD och firmatecknare kan fortsätta."

\subsubsection{Teknisk sammanfattning}

Riskfrågor Steg 1 består av ett enda POST endpoint som tar emot 8 grundfrågor och gör omfattande validering mot externa API:er.

\textbf{Externa anrop (total):}
\begin{itemize}[noitemsep]
  \item 1× Bolagsverket (organisationer) — Validera org.nr + VD
  \item 1× SPAR (personnummer) — Validera VD personnummer
  \item 1-2× Roaring (beneficial-owners) — Hämta VH (primary + optional fallback)
  \item N× Roaring (PEP) — Kontrollera VD + VH mot PEP-listor (N = antal VH)
  \item N× Roaring (Sanctions) — Kontrollera VD + VH mot sanktionslistor
  \item N× Roaring (BusinessProhibition) — Kontrollera VD + VH mot näringsförbud
\end{itemize}

\textbf{Svarstid:} 3-5 sekunder beroende på antal VH och nätverkshastighet

\textbf{Databas-operationer:}
\begin{itemize}[noitemsep]
  \item 1× SELECT (verifiera ägarskap av onboardingId)
  \item 1× UPDATE (spara riskfrågor + triggers + externa API-data)
  \item 3-9× INSERT (logga alla externa API-anrop i external\_api\_responses)
\end{itemize}

\textbf{Conditional routing:}
\begin{itemize}[noitemsep]
  \item OM alla triggers = false → Nästa steg = /identitetskontroll (hoppar över Steg 2)
  \item OM minst en trigger = true → Nästa steg = /riskfragor/steg2
\end{itemize}

\textbf{Frontend uppdatering:}
\begin{itemize}[noitemsep]
  \item Spara triggers i localStorage
  \item Spara bolagsverketData och sparData i localStorage (för visning)
  \item Uppdatera currentStep till 2
  \item Navigera till nextStep URL
\end{itemize}

\textbf{Blockerande fel:}
\begin{itemize}[noitemsep]
  \item Personnummer matchar inte VD/firmatecknare → 403 Forbidden
  \item Organisationsnummer finns inte → 404 Not Found
  \item Person är sanktionerad → 403 Forbidden (REJECT omedelbart)
  \item Person har näringsförbud → 403 Forbidden (REJECT omedelbart)
\end{itemize}

\textbf{Icke-blockerande varningar:}
\begin{itemize}[noitemsep]
  \item Person är PEP → Trigger showPEP = true (kräver fördjupning i Steg 2)
  \item Person ej folkbokförd → Varning till compliance-team (men fortsätt onboarding)
  \item Komplex ägarstruktur (totalCoverage < 50\%) → Flagga för manuell granskning
\end{itemize}


\newpage

% ========================================================================
% SEKTION 2B: RISKFRÅGOR STEG 2 (FÖRDJUPADE FRÅGOR)
% ========================================================================

\section{SEKTION 2B: Riskfrågor – Steg 2 (Fördjupade frågor)}

\subsubsection{POST /api/onboarding/\{onboardingId\}/riskfragor/steg2}

\textbf{Beskrivning:}\\
Sparar fördjupade riskfrågor baserat på triggers från Steg 1. Endast frågor där trigger = true visas för användaren. Uppdaterar riskpoäng och beräknar om enhanced due diligence (EDD) krävs.

\textbf{Request Headers:}
\begin{lstlisting}[language=json]
{
  "Authorization": "Bearer <access_token>",
  "Content-Type": "application/json"
}
\end{lstlisting}

\textbf{Request Body (Exempel - Alla triggers = true):}
\begin{lstlisting}[language=json]
{
  "utlandskaTransaktioner": {
    "lander": "Norge, Danmark, Tyskland",
    "andel": "20%",
    "syfte": "Konsultuppdrag för nordiska kunder",
    "valutor": "NOK, DKK, EUR"
  },
  "kontanthantering": {
    "omfattning": "Sällan",
    "maxBelopp": "5000 SEK",
    "rutiner": "Kvitton sparas, bokförs samma dag"
  },
  "verksamhetsandring": {
    "vad": "Lagt till online-kurser som ny intäktskälla",
    "nar": "2024-06",
    "varfor": "Efterfrågan från kunder"
  },
  "pepDetaljer": {
    "befattning": "Före detta kommunalråd",
    "period": "2015-2019",
    "familjerelation": "Nej"
  }
}
\end{lstlisting}

\textbf{Request Body (Exempel - Endast utländska = true):}
\begin{lstlisting}[language=json]
{
  "utlandskaTransaktioner": {
    "lander": "Norge, Danmark",
    "andel": "10%",
    "syfte": "Konsultuppdrag",
    "valutor": "NOK, DKK"
  }
}
\end{lstlisting}

\textbf{Response (200 OK - Standard Risk):}
\begin{lstlisting}[language=json]
{
  "success": true,
  "riskScore": 35,
  "riskLevel": "STANDARD",
  "eddRequired": false,
  "message": "Standardrisk - ingen fördjupad granskning krävs. Du kan gå vidare till identitetskontroll.",
  "nextStep": "/identitetskontroll",
  "riskFactors": [
    {
      "factor": "Utländska transaktioner",
      "severity": "LOW",
      "points": 15,
      "reason": "Nordiska länder (lågriskländer), <25% av omsättning"
    },
    {
      "factor": "Kontanthantering",
      "severity": "LOW",
      "points": 10,
      "reason": "Sällan, små belopp, goda rutiner"
    },
    {
      "factor": "Verksamhetsändring",
      "severity": "LOW",
      "points": 10,
      "reason": "Naturlig expansion, logisk förklaring"
    }
  ]
}
\end{lstlisting}

\textbf{Response (200 OK - Enhanced Due Diligence Required):}
\begin{lstlisting}[language=json]
{
  "success": true,
  "riskScore": 85,
  "riskLevel": "HIGH",
  "eddRequired": true,
  "message": "Förhöjd risk - fördjupad due diligence krävs. Vi behöver granska fler dokument.",
  "nextStep": "/identitetskontroll",
  "riskFactors": [
    {
      "factor": "PEP-status",
      "severity": "HIGH",
      "points": 40,
      "reason": "Politiskt exponerad person (före detta kommunalråd)"
    },
    {
      "factor": "Utländska transaktioner",
      "severity": "MEDIUM",
      "points": 25,
      "reason": "Flera länder, >10% av omsättning"
    },
    {
      "factor": "Kontanthantering",
      "severity": "MEDIUM",
      "points": 20,
      "reason": "Privatpersoner som kunder, kontanter förekommer"
    }
  ],
  "eddRequirements": [
    "Specificera källa till förmögenhet för PEP",
    "Tillhandahåll kontrakt för utländska uppdrag",
    "Dokumentera kontanthanteringsrutiner skriftligt",
    "Förklara relation till utländska partners"
  ]
}
\end{lstlisting}

\textbf{Response (400 Bad Request):}
\begin{lstlisting}[language=json]
{
  "success": false,
  "error": "Obligatoriska fält saknas",
  "code": "MISSING_REQUIRED_FIELDS",
  "details": {
    "utlandskaTransaktioner": "Detta fält är obligatoriskt när showUtlandska = true",
    "pepDetaljer": "Detta fält är obligatoriskt när showPEP = true"
  }
}
\end{lstlisting}

\textbf{Response (403 Forbidden):}
\begin{lstlisting}[language=json]
{
  "success": false,
  "error": "Onboarding-process ägs inte av inloggad användare",
  "code": "UNAUTHORIZED_ACCESS"
}
\end{lstlisting}

\textbf{Business Logic:}
\begin{enumerate}[noitemsep]
  \item \textbf{Hämta user\_id från JWT och verifiera ägarskap:}
    \begin{itemize}[noitemsep]
      \item SELECT från \texttt{onboarding\_processes} WHERE id = \texttt{onboardingId} AND user\_id = \texttt{user\_id\_from\_jwt}
      \item Hämta \texttt{triggers} från Steg 1 (vilka frågor ska ställas)
    \end{itemize}
  \item \textbf{Validera att endast relevanta fält är ifyllda:}
    \begin{itemize}[noitemsep]
      \item Om triggers.showUtlandska = false → utlandskaTransaktioner får INTE finnas
      \item Om triggers.showUtlandska = true → utlandskaTransaktioner MÅSTE finnas
      \item Samma logik för kontanthantering, verksamhetsandring, pepDetaljer
    \end{itemize}
  \item \textbf{Beräkna riskpoäng (0-100):}
    \begin{itemize}[noitemsep]
      \item \textbf{Utländska transaktioner:}
        \begin{itemize}[noitemsep]
          \item Lågriskländer (Norden, EU15) + <25\% andel → 10-20 poäng (LOW)
          \item Medelriskländer (EU, OECD) + 25-50\% andel → 20-35 poäng (MEDIUM)
          \item Högriskländer (FATF graylisting) + >50\% andel → 40-60 poäng (HIGH)
        \end{itemize}
      \item \textbf{Kontanthantering:}
        \begin{itemize}[noitemsep]
          \item Sällan + <10,000 SEK + goda rutiner → 5-10 poäng (LOW)
          \item Ibland + 10,000-50,000 SEK + acceptabla rutiner → 15-25 poäng (MEDIUM)
          \item Ofta + >50,000 SEK + bristfälliga rutiner → 30-50 poäng (HIGH)
        \end{itemize}
      \item \textbf{Verksamhetsändring:}
        \begin{itemize}[noitemsep]
          \item Naturlig expansion + logisk förklaring → 5-10 poäng (LOW)
          \item Stor förändring + vag förklaring → 15-25 poäng (MEDIUM)
          \item Komplett pivoting + ingen förklaring → 30-40 poäng (HIGH)
        \end{itemize}
      \item \textbf{PEP-status:}
        \begin{itemize}[noitemsep]
          \item Före detta (>2 år sedan) + lokal nivå → 30-40 poäng (HIGH)
          \item Nuvarande + lokal nivå → 50-60 poäng (HIGH)
          \item Nuvarande + nationell/internationell nivå → 70-90 poäng (CRITICAL)
          \item Familj till PEP → +20 poäng extra
        \end{itemize}
    \end{itemize}
  \item \textbf{Avgör risk level:}
    \begin{itemize}[noitemsep]
      \item 0-30 poäng → STANDARD (eddRequired = false)
      \item 31-60 poäng → MEDIUM (eddRequired = true, mild EDD)
      \item 61-100 poäng → HIGH (eddRequired = true, omfattande EDD)
    \end{itemize}
  \item \textbf{Generera EDD-krav baserat på riskfaktorer}
  \item \textbf{Uppdatera onboarding\_processes:}
    \begin{itemize}[noitemsep]
      \item Spara riskfragor\_steg2, risk\_score, risk\_level, edd\_required
      \item Uppdatera current\_step till 3 (identitetskontroll)
    \end{itemize}
\end{enumerate}

\textbf{Databas:}
\begin{lstlisting}[language=sql]
-- Hamta triggers fran Steg 1 for validering
SELECT 
  triggers,
  riskfragor_steg1
FROM onboarding_processes
WHERE id = $1                    -- onboardingId (UUID)
  AND user_id = $2;              -- Fran JWT (security check)

-- Uppdatera med Steg 2 svar och riskbedömning
UPDATE onboarding_processes
SET 
  riskfragor_steg2 = $1::jsonb,
  risk_score = $2,               -- 0-100
  risk_level = $3,               -- STANDARD/MEDIUM/HIGH
  edd_required = $4,             -- boolean
  risk_factors = $5::jsonb,      -- Array av riskfaktorer
  current_step = 3,
  status = 'in_progress',
  updated_at = NOW()
WHERE id = $6                    -- onboardingId (UUID)
  AND user_id = $7;              -- Fran JWT token (security check)
\end{lstlisting}

\textbf{Database Schema (behövs i Loop 4):}
\begin{lstlisting}[language=sql]
-- Lagg till riskbedömning kolumner i onboarding_processes
ALTER TABLE onboarding_processes
ADD COLUMN IF NOT EXISTS riskfragor_steg2 JSONB,
ADD COLUMN IF NOT EXISTS risk_score INTEGER DEFAULT 0 CHECK (risk_score >= 0 AND risk_score <= 100),
ADD COLUMN IF NOT EXISTS risk_level VARCHAR(20) CHECK (risk_level IN ('STANDARD', 'MEDIUM', 'HIGH')),
ADD COLUMN IF NOT EXISTS edd_required BOOLEAN DEFAULT FALSE,
ADD COLUMN IF NOT EXISTS risk_factors JSONB DEFAULT '[]'::jsonb;

-- Index for risk-based queries
CREATE INDEX IF NOT EXISTS idx_onboarding_risk_level 
  ON onboarding_processes (risk_level);
CREATE INDEX IF NOT EXISTS idx_onboarding_edd_required 
  ON onboarding_processes (edd_required);
\end{lstlisting}

\textbf{PTL-relevans:}
\begin{itemize}[noitemsep]
  \item \textbf{PTL 2 kap 2 §:} Riskbaserad bedömning krävs
  \item \textbf{PTL 2 kap 8 §:} Förhöjd risk vid:
    \begin{itemize}[noitemsep]
      \item PEP (Politically Exposed Person)
      \item Högriskländer enligt FATF
      \item Ovanliga transaktionsmönster
      \item Omfattande kontanthantering
    \end{itemize}
  \item \textbf{PTL 3 kap 12-14 §§:} Enhanced Due Diligence (EDD) för högrisk
  \item \textbf{PTL 4 kap 3 §:} Transaktionsövervakning
\end{itemize}

\newpage

\section{SEKTION 3: Identitetskontroll}

\textit{TODO: Dokumentera identitetskontroll och dokumentation}

\section{SEKTION 4: Jämförelse mot externa register}

\textit{TODO: Dokumentera resultatslides från Bolagsverket/Roaring.io}

\section{SEKTION 5: Företagsdokumentation}

\textit{TODO: Dokumentera dokumentuppladdning}

\section{SEKTION 6: Ekonomisk analys}

\textit{TODO: Dokumentera SIE-fil, IBAN, Klarna Kosma integration}

\section{SEKTION 7: Djupgranskning och avtal}

\textit{TODO: Dokumentera bokföringsanalys, riskbedömning, BankID-signering}

\end{document}
